\documentclass[11pt,a4paper]{article}
\usepackage[utf8]{inputenc}
\usepackage[T1]{fontenc}
\usepackage[french]{babel}
\usepackage[a4paper,margin=2.2cm]{geometry}
\usepackage{amsmath,amssymb,booktabs,tabularx,array,hyperref}
\hypersetup{colorlinks=true,linkcolor=blue,urlcolor=blue}
\newcolumntype{Y}{>{\raggedright\arraybackslash}X}

\title{Rapport de verrouillage scientifique\\Phase 1/10}
\author{Projet de stéganographie adaptative par parcours de graphe}
\date{12 juin 2026}

\begin{document}
\maketitle

\begin{abstract}
Cette phase fixe l'objet transmis, le modèle de menace, les questions de
recherche, les revendications d'originalité admissibles et les critères
d'abandon. Le résultat principal est un recentrage : le stego est une trace
de parcours temporel valide, et non un graphe modifié. L'originalité ne peut
pas reposer sur l'usage isolé d'un chemin, d'un RNN ou de la logique floue,
car ces composants possèdent des antériorités. La contribution visée est leur
intégration vérifiable dans un canal temporel appris et évalué contre des
stéganalystes externes sur des trajectoires ou sessions réelles.
\end{abstract}

\section{Décision scientifique}

Soit un graphe temporel observé jusqu'à l'instant \(T\). Une trace admissible
est définie par
\[
\tau=(v_0,e_1,t_1,v_1,\ldots,e_L,t_L,v_L),\qquad
t_1<\cdots<t_L\leq T.
\]
La méthode n'ajoute ni sommet ni arête et ne modifie pas les timestamps. Le
message chiffré détermine les choix successifs parmi les continuations
admissibles. Le canal est donc une stéganographie générative ou sélective de
parcours.

Un modèle causal estime
\[
q_\theta(a_i\mid\tau_{<i},G_{\leq t_i}).
\]
Un codeur entropique associe les bits aux continuations. Un contrôleur flou
réduit le débit ou impose l'abstention selon l'entropie, la calibration, le
risque de détection, la charge restante et la fragilité du canal.

\section{Extension formelle de la méthode}

Quatre distributions sont séparées : \(P^\star\), distribution réelle
inconnue; \(Q_\theta\), modèle causal appris; \(C_\psi\), distribution après
contrôle; et \(S\), distribution induite par le codeur. La politique
\(\pi=(\theta,\psi,\phi)\) optimise
\[
\mathcal J(\pi)=R(\pi)
-\lambda_dD(S_\pi,P^\star)
-\lambda_r\operatorname{Risk}_{Eve}(\pi)
-\lambda_e\operatorname{BER}(\pi)
-\lambda_c\operatorname{Cost}(\pi).
\]
Le point publié est toutefois choisi sur validation par maximisation du débit
sous contraintes d'AUC, de BER et d'utilité.

Le coût local d'une continuation est
\[
c_i(a)=
\alpha[-\log Q_\theta(a\mid h_i)]
+\beta r_{\rm steg}(a,h_i)
+\gamma r_{\rm canal}(a,h_i)
+\xi r_{\rm cal}(a,h_i).
\]
Le signal du stéganalyste modifie donc les probabilités disponibles avant
l'encodage. Le contrôleur Takagi--Sugeno choisit un budget de bits et l'un des
quatre modes \(\mathtt{EMBED}\), \(\mathtt{COVER}\), \(\mathtt{PAUSE}\) ou
\(\mathtt{STOP}\).

La divergence KL n'étant pas triangulaire, la décomposition théorique utilise
des hypothèses de rapport de vraisemblance. Si
\[
\log\frac{Q_\theta}{P^\star}\leq\varepsilon_{\rm mod},\qquad
\log\frac{C_\psi}{Q_\theta}\leq\varepsilon_{\rm ctrl},\qquad
D_{\rm KL}(S\Vert C_\psi)\leq\varepsilon_{\rm code},
\]
sur le support stego, alors
\[
D_{\rm KL}(S\Vert P^\star)
\leq\varepsilon_{\rm mod}
+\varepsilon_{\rm ctrl}
+\varepsilon_{\rm code}.
\]
Cette proposition structure l'analyse sans prétendre fournir une sécurité
parfaite contre une distribution réelle inconnue.

\section{Modèle de menace}

Alice et Bob partagent une clé, une version du modèle, un graphe de référence
et une règle déterministe de construction des candidats. Eve connaît
l'algorithme, l'architecture et les données publiques, mais pas la clé. Elle
peut entraîner des détecteurs différents de celui utilisé pendant
l'apprentissage.

Le scénario principal est passif. Le scénario secondaire autorise suppression,
insertion, réordonnancement, troncature et bruit temporel. La correction
requise dans le canal passif est
\[
\Pr[\operatorname{Dec}_k(\operatorname{Enc}_k(m,G),G)=m]=1.
\]

La sécurité parfaite vis-à-vis de la distribution réelle n'est pas
revendiquée. Il faut distinguer la fidélité à \(q_\theta\), mesurable
conditionnellement au modèle, de la détectabilité face à la distribution
réelle inconnue \(P^\star\).

\section{Contrôle d'antériorité}

\begingroup
\footnotesize
\setlength{\tabcolsep}{3pt}
\begin{tabularx}{0.96\textwidth}{@{}YYYYY@{}}
\toprule
Travail & Chemin & Distribution & Adaptation & Limite pour notre projet \\
\midrule
BIND/AdaBIND & Non & Statique & Ajout d'arêtes & Pas de dynamique ni
stéganalyse apprise \\
Graph-Stega & Oui & Graphe de connaissances & Non & Invalide la primauté
générale sur les chemins \\
ADG/iMEC & Non & Autorégressive & Variable & Références fortes pour
l'encodage distributionnel \\
ADLM/FreStega & Non & Modèle de langue & Entropie/domaine & Invalident la
primauté sur l'adaptation entropique \\
Joshi--Bhand & Non & Image statique & Flou Mamdani & Invalide la primauté
générale sur la logique floue \\
Mesh FPD & Non & Coût stéganalytique & Distorsion/STC & Baseline
méthodologique d'adaptation au détecteur \\
\bottomrule
\end{tabularx}
\endgroup

Les formulations ``première stéganographie par chemin'', ``première
stéganographie floue'' et ``sécurité parfaite'' sont interdites. La
revendication provisoire porte sur la combinaison spécifique d'un parcours
temporel causal, d'un encodage distributionnel, d'un contrôle flou avec
abstention et d'une stéganalyse externe.

\section{Questions et hypothèses}

\begin{enumerate}
\item Le modèle temporel est-il mieux calibré que les parcours uniforme,
DeepWalk et node2vec?
\item Le codeur conserve-t-il la distribution du modèle sans erreur de
décodage dans le canal passif?
\item Le contrôleur flou améliore-t-il le front
capacité--détectabilité--fiabilité face à un seuil fixe et une MLP?
\item Le gain subsiste-t-il contre des détecteurs non vus et en transfert?
\item Quel niveau de redondance protège le message sous altération du canal?
\end{enumerate}

Le contrôleur flou doit améliorer l'hypervolume d'au moins 5\,\% avec un
intervalle bootstrap excluant zéro. Au point d'exploitation, l'AUC externe
doit être au plus de 0,60 sur trois domaines, le BER passif doit être nul et
le BER corrigé doit être inférieur ou égal à \(10^{-3}\) au niveau d'attaque
retenu.

\section{Politique des données}

Les données sont séparées en deux niveaux. Le niveau A contient des
trajectoires ou sessions réellement observées et fonde les conclusions de
furtivité. Le niveau B contient des flux d'interactions temporelles, notamment
TGB, utiles pour l'apprentissage et le passage à l'échelle mais insuffisants
pour qualifier seuls un parcours de comportement naturel.

Porto Taxi, T-Drive, GeoLife et des sessions réelles de recommandation sont
des candidats de niveau A. Leur admission exige une licence vérifiée, une
provenance, un checksum et une segmentation reproductible. Aucune donnée
synthétique ne figurera dans les résultats scientifiques.

\section{Critères de pivot}

\begin{itemize}
\item Retirer la logique floue si elle ne bat pas le seuil fixe.
\item Retirer GRU ou TGN si une architecture plus simple est aussi bien
calibrée et aussi furtive.
\item Optimiser l'alignement au domaine si l'erreur du modèle domine l'erreur
d'encodage.
\item Réserver TGB au préentraînement si la sémantique des parcours n'est pas
observable.
\item Abandonner le canal si l'AUC reste supérieure à 0,65 à toute charge
utile exploitable sur trois datasets.
\end{itemize}

\section{Verdict}

Le sujet est scientifiquement prometteur, mais sa valeur Q1 dépendra d'une
démonstration causale de l'apport de chaque composant. Une simple combinaison
GNN--GRU--flou ne constitue pas une contribution suffisante. La phase est
gelée provisoirement; les revendications de primauté restent suspendues
jusqu'à la clôture des recherches dans les bases institutionnelles.

\end{document}
